\documentclass[UTF8,fontset=windows,11pt]{ctexart}
\usepackage[a4paper,margin=1.85cm,headheight=24pt,footskip=24pt]{geometry}
\usepackage{xcolor}
\usepackage{graphicx}
\usepackage{enumitem}
\usepackage{tcolorbox}
\usepackage{fancyhdr}
\usepackage{lastpage}
\usepackage{titlesec}
\usepackage{hyperref}
\usepackage{tabularx}
\usepackage{booktabs}
\tcbuselibrary{skins,breakable}

\definecolor{examnavy}{HTML}{1E3A5F}
\definecolor{examblue}{HTML}{2563A7}
\definecolor{examred}{HTML}{C2410C}
\definecolor{examgray}{HTML}{4B5563}
\definecolor{examline}{HTML}{D7DEE8}
\definecolor{exambg}{HTML}{F8FAFC}
\definecolor{choicepink}{HTML}{DB2777}
\definecolor{answergreen}{HTML}{15803D}
\definecolor{analysispurple}{HTML}{7C3AED}
\definecolor{analysisbg}{HTML}{F5F3FF}

\hypersetup{colorlinks=true,linkcolor=examnavy,urlcolor=examblue}
\setlength{\parindent}{0pt}
\setlength{\parskip}{0.28em}
\linespread{1.08}

\pagestyle{fancy}
\fancyhf{}
\lhead{\small\textcolor{examnavy}{施工章节测试}}
\rhead{\small\textcolor{examred}{复习指南对应题 答案解析版}}
\cfoot{\small\textcolor{examgray}{第 \thepage\ 页 / 共 \pageref{LastPage} 页}}
\renewcommand{\headrulewidth}{0.4pt}
\renewcommand{\headrule}{\hbox to\headwidth{\color{examline}\leaders\hrule height \headrulewidth\hfill}}

\titleformat{\section}
  {\Large\bfseries\color{examnavy}}
  {\thesection}
  {0.8em}
  {}
\titlespacing*{\section}{0pt}{1.1em}{0.55em}

\newcommand{\chaptermeta}[1]{%
  \vspace{-0.25em}
  {\small\textcolor{examgray}{#1}}\par
  \vspace{0.35em}
  {\color{examline}\rule{\linewidth}{0.55pt}}\vspace{0.25em}
}

\newtcolorbox{analysisbox}[2]{
  enhanced,
  breakable,
  colback=white,
  colframe=examline,
  boxrule=0.55pt,
  arc=1mm,
  left=3mm,
  right=3mm,
  top=2.4mm,
  bottom=2mm,
  before skip=0.65em,
  after skip=0.72em,
  borderline west={1.8pt}{0pt}{analysispurple},
  title={\textcolor{examnavy}{#1}\hfill\textcolor{examgray}{#2}},
  coltitle=examnavy,
  colbacktitle=exambg,
  fonttitle=\bfseries,
  boxed title style={
    colback=exambg,
    colframe=examline,
    boxrule=0.35pt,
    arc=0.8mm
  },
  attach boxed title to top left={xshift=2mm,yshift=-2mm},
  top=6mm
}

\newcommand{\inlineanswer}[1]{%
  \vspace{0.25em}
  \begin{tcolorbox}[
    enhanced,
    breakable,
    colback=answergreen!6,
    colframe=answergreen,
    boxrule=0.55pt,
    arc=0.8mm,
    left=2.5mm,
    right=2.5mm,
    top=1.2mm,
    bottom=1.2mm,
    borderline west={1.6pt}{0pt}{answergreen}
  ]
  \textbf{\textcolor{answergreen}{答案：}}#1
  \end{tcolorbox}
}

\newcommand{\analysiscontent}[1]{%
  \vspace{0.15em}
  \begin{tcolorbox}[
    enhanced,
    breakable,
    colback=analysisbg,
    colframe=analysispurple!30,
    boxrule=0.45pt,
    arc=1mm,
    left=2.5mm,
    right=2.5mm,
    top=1mm,
    bottom=1mm
  ]
  {\small\color{examgray}#1}
  \end{tcolorbox}
}

\begin{document}



\begin{titlepage}
\vspace*{1.2cm}
{\Huge\bfseries\textcolor{examnavy}{施工章节测试}}\par
\vspace{0.35em}
{\LARGE\bfseries\textcolor{examred}{复习指南对应题 答案解析版}}\par
\vspace{0.8em}
{\color{examline}\rule{0.82\linewidth}{0.8pt}}\par
\vspace{1.1em}
{\large 题目 + 答案速查表 + 每题详细解析；便于刷题后对照复盘。}\par
\vspace{2.2em}
\begin{tcolorbox}[
  enhanced,
  width=0.78\linewidth,
  colback=exambg,
  colframe=examline,
  boxrule=0.6pt,
  arc=1mm,
  left=4mm,
  right=4mm,
  top=3mm,
  bottom=3mm
]
\Large
\textcolor{examnavy}{章节数}\quad\textbf{11}\qquad
\textcolor{examnavy}{题目数}\quad\textbf{145}\par
\vspace{0.45em}
\normalsize\textcolor{examgray}{建议先独立刷题，再对照答案与解析复盘。}
\end{tcolorbox}
\vfill
{\small\textcolor{examgray}{依据施工章节测试复习指南对应题刷题版整理}}
\end{titlepage}
\tableofcontents
\newpage


\section{答案速查表}

\chaptermeta{答题完毕后，对照此表快速核对答案。}

\noindent\textbf{\textcolor{examnavy}{第一章  土方工程}}\\\vspace{0.3em}

\begin{tabularx}{\linewidth}{|c|X|X|X|X|X|X|X|X|X|X|}
\hline
\textbf{题号} & \textbf{1} & \textbf{2} & \textbf{3} & \textbf{4} & \textbf{5} & \textbf{6} & \textbf{7} & \textbf{8} & \textbf{9} & \textbf{10} \\
\hline
Q1-10 & A & B & C & D & C & A & ACDE & D & ABCD & A \\
\hline
Q11-20 & D & B & C & BCE & C & B & ABC & D & ACD & ACE \\
\hline
Q21-30 & B & A & B & BD & BCDE & A & BCE & B & A & C \\
\hline
\end{tabularx}
\vspace{0.6em}

\noindent\textbf{\textcolor{examnavy}{第二章  桩基础工程}}\\\vspace{0.3em}

\begin{tabularx}{\linewidth}{|c|X|X|X|X|X|X|X|X|X|X|}
\hline
\textbf{题号} & \textbf{1} & \textbf{2} & \textbf{3} & \textbf{4} & \textbf{5} & \textbf{6} & \textbf{7} & \textbf{8} & \textbf{9} & \textbf{10} \\
\hline
Q31-37 & D & D & BE & CE & D & A & B &  &  &  \\
\hline
\end{tabularx}
\vspace{0.6em}

\noindent\textbf{\textcolor{examnavy}{第三章  基坑工程}}\\\vspace{0.3em}

\begin{tabularx}{\linewidth}{|c|X|X|X|X|X|X|X|X|X|X|}
\hline
\textbf{题号} & \textbf{1} & \textbf{2} & \textbf{3} & \textbf{4} & \textbf{5} & \textbf{6} & \textbf{7} & \textbf{8} & \textbf{9} & \textbf{10} \\
\hline
Q38-47 & C & D & ABD & AB & A & C & BE & AE & A & A \\
\hline
\end{tabularx}
\vspace{0.6em}

\noindent\textbf{\textcolor{examnavy}{第四章  混凝土结构工程}}\\\vspace{0.3em}

\begin{tabularx}{\linewidth}{|c|X|X|X|X|X|X|X|X|X|X|}
\hline
\textbf{题号} & \textbf{1} & \textbf{2} & \textbf{3} & \textbf{4} & \textbf{5} & \textbf{6} & \textbf{7} & \textbf{8} & \textbf{9} & \textbf{10} \\
\hline
Q48-57 & A & C & B & BD & B & ABC & C & D & B & D \\
\hline
Q58-62 & BCD & ADE & C & A & CDE &  &  &  &  &  \\
\hline
\end{tabularx}
\vspace{0.6em}

\noindent\textbf{\textcolor{examnavy}{第五章  预应力混凝土工程}}\\\vspace{0.3em}

\begin{tabularx}{\linewidth}{|c|X|X|X|X|X|X|X|X|X|X|}
\hline
\textbf{题号} & \textbf{1} & \textbf{2} & \textbf{3} & \textbf{4} & \textbf{5} & \textbf{6} & \textbf{7} & \textbf{8} & \textbf{9} & \textbf{10} \\
\hline
Q63-69 & A & C & D & A & C & D & ACE &  &  &  \\
\hline
\end{tabularx}
\vspace{0.6em}

\noindent\textbf{\textcolor{examnavy}{第六章  砌筑工程}}\\\vspace{0.3em}

\begin{tabularx}{\linewidth}{|c|X|X|X|X|X|X|X|X|X|X|}
\hline
\textbf{题号} & \textbf{1} & \textbf{2} & \textbf{3} & \textbf{4} & \textbf{5} & \textbf{6} & \textbf{7} & \textbf{8} & \textbf{9} & \textbf{10} \\
\hline
Q70-79 & D & C & A & D & C & D & A & C & ABC & D \\
\hline
Q80-89 & C & BCD & B & C & A & B & C & D & C & C \\
\hline
Q90-94 & D & BE & C & B & AD &  &  &  &  &  \\
\hline
\end{tabularx}
\vspace{0.6em}

\noindent\textbf{\textcolor{examnavy}{第七章  脚手架工程}}\\\vspace{0.3em}

\begin{tabularx}{\linewidth}{|c|X|X|X|X|X|X|X|X|X|X|}
\hline
\textbf{题号} & \textbf{1} & \textbf{2} & \textbf{3} & \textbf{4} & \textbf{5} & \textbf{6} & \textbf{7} & \textbf{8} & \textbf{9} & \textbf{10} \\
\hline
Q95-99 & BC & CE & B & AB & AE &  &  &  &  &  \\
\hline
\end{tabularx}
\vspace{0.6em}

\noindent\textbf{\textcolor{examnavy}{第八章  装饰工程(抹灰工程)}}\\\vspace{0.3em}

\begin{tabularx}{\linewidth}{|c|X|X|X|X|X|X|X|X|X|X|}
\hline
\textbf{题号} & \textbf{1} & \textbf{2} & \textbf{3} & \textbf{4} & \textbf{5} & \textbf{6} & \textbf{7} & \textbf{8} & \textbf{9} & \textbf{10} \\
\hline
Q100-101 & BDE & C &  &  &  &  &  &  &  &  \\
\hline
\end{tabularx}
\vspace{0.6em}

\noindent\textbf{\textcolor{examnavy}{第九章  防水工程}}\\\vspace{0.3em}

\begin{tabularx}{\linewidth}{|c|X|X|X|X|X|X|X|X|X|X|}
\hline
\textbf{题号} & \textbf{1} & \textbf{2} & \textbf{3} & \textbf{4} & \textbf{5} & \textbf{6} & \textbf{7} & \textbf{8} & \textbf{9} & \textbf{10} \\
\hline
Q102-108 & A & BDE & BE & ABCE & A & ACE & B &  &  &  \\
\hline
\end{tabularx}
\vspace{0.6em}

\noindent\textbf{\textcolor{examnavy}{第十章  流水施工原理}}\\\vspace{0.3em}

\begin{tabularx}{\linewidth}{|c|X|X|X|X|X|X|X|X|X|X|}
\hline
\textbf{题号} & \textbf{1} & \textbf{2} & \textbf{3} & \textbf{4} & \textbf{5} & \textbf{6} & \textbf{7} & \textbf{8} & \textbf{9} & \textbf{10} \\
\hline
Q109-118 & C & C & A & B & B & B & C & D & B & C \\
\hline
Q119-122 & D & C & B & C &  &  &  &  &  &  \\
\hline
\end{tabularx}
\vspace{0.6em}

\noindent\textbf{\textcolor{examnavy}{第十一章  网络计划技术}}\\\vspace{0.3em}

\begin{tabularx}{\linewidth}{|c|X|X|X|X|X|X|X|X|X|X|}
\hline
\textbf{题号} & \textbf{1} & \textbf{2} & \textbf{3} & \textbf{4} & \textbf{5} & \textbf{6} & \textbf{7} & \textbf{8} & \textbf{9} & \textbf{10} \\
\hline
Q123-132 & D & C & B & A & A & D & A & A & B & B \\
\hline
Q133-142 & D & BC & B & C & D & C & D & B & BD & AB \\
\hline
Q143-145 & BCD & A & BD &  &  &  &  &  &  &  \\
\hline
\end{tabularx}
\vspace{0.6em}

\newpage

\section{答案解析}

\chaptermeta{以下对每道题给出正确答案解析和易错点提示。}

\subsection{第一章  土方工程}

\begin{analysisbox}{第 1 题}{答案：A}

\inlineanswer{A}

\analysiscontent{填土压实质量控制指标为土的干密度(又称干容重),即单位体积内土的固体颗粒质量。密实度反映压实程度而非直接控制指标;压缩比和可松性均不是压实控制指标。}

\end{analysisbox}

\begin{analysisbox}{第 2 题}{答案：B}

\inlineanswer{B}

\analysiscontent{含水量 w = 水的质量/干土质量 x 100\%,即水的质量与干土质量之比的百分数。选项A混淆为湿土质量;C、D分别描述了饱和度和含水比的概念。}

\end{analysisbox}

\begin{analysisbox}{第 3 题}{答案：C}

\inlineanswer{C}

\analysiscontent{自然状态挖方量 V1 = 1000m3,可松性系数 Ks = 1.25,松散状态下需运走的土方量 V2 = Ks * V1 = 1.25 x 1000 = 1250m3。最终可松系数 Ks=1.05 用于回填计算,与运输量无关。}

\end{analysisbox}

\begin{analysisbox}{第 4 题}{答案：D}

\inlineanswer{D}

\analysiscontent{推土机操作灵活、行驶速度快、转移方便,可爬30度左右的缓坡,集切土、推土和卸土功能于一体,所需工作面小。正铲/反铲挖掘机不具备推土功能;铲运机不善于爬陡坡。}

\end{analysisbox}

\begin{analysisbox}{第 5 题}{答案：C}

\inlineanswer{C}

\analysiscontent{铲运机适用于大面积场地平整和填筑路基堤坝,但在沼泽区施工时地基承载力不足,机械容易下陷,故不适宜。含水量不高的松土是铲运机的适用工况。}

\end{analysisbox}

\begin{analysisbox}{第 6 题}{答案：A}

\inlineanswer{A}

\analysiscontent{正铲挖掘机的挖土特点是前进向上,强制切土。挖掘机向前进方向推进,铲斗自下而上切削土体。B为反铲特点;C为拉铲特点;D为抓铲特点。}

\end{analysisbox}

\begin{analysisbox}{第 7 题}{答案：ACDE}

\inlineanswer{ACDE}

\analysiscontent{土方机械选择应综合考虑:土方工程的范围大小(A)、不同施工机械的组合(C)、土的含水量(D)、施工合同工期要求(E)。施工合同的要求(B)属于商务条款,不属于技术选型依据。}

\end{analysisbox}

\begin{analysisbox}{第 8 题}{答案：D}

\inlineanswer{D}

\analysiscontent{可松性是土方平衡调配的关键参数:挖出的自然土变成松散土,回填压实后体积又发生变化,直接决定挖填平衡、弃土量和借土量。天然含水量、天然密度、密实度均不是平衡调配的核心参数。}

\end{analysisbox}

\begin{analysisbox}{第 9 题}{答案：ABCD}

\inlineanswer{ABCD}

\analysiscontent{内摩擦角φ是土的抗剪强度指标之一(A正确),反映了土颗粒之间的摩擦特性(B正确),力学上等价于块体在斜面上的临界自稳角(C正确),可用于分析边坡稳定性(D正确)。E错误:在内摩擦角度内块体是稳定的,超过此角才会滑动,选项说反了。}

\end{analysisbox}

\begin{analysisbox}{第 10 题}{答案：A}

\inlineanswer{A}

\analysiscontent{放坡挖土是最简单、最经济的基坑开挖方案,仅在地质条件和场地条件许可、开挖深度不大时采用。中心岛、盆式、逆作法均属于有支护结构的开挖方案。}

\end{analysisbox}

\begin{analysisbox}{第 11 题}{答案：D}

\inlineanswer{D}

\analysiscontent{放坡开挖是唯一不需要支护结构的挖土方案。中心岛式、盆式、逆作法均需要围护结构配合。}

\end{analysisbox}

\begin{analysisbox}{第 12 题}{答案：B}

\inlineanswer{B}

\analysiscontent{中心岛式挖土:先开挖基坑周边土体,中间保留土墩作为中心岛,可作为临时施工场地和运输通道(B正确)。A:四边应开挖而非留土坡;C:有利于支护变形控制描述不准确;D:多用于有支护工程。}

\end{analysisbox}

\begin{analysisbox}{第 13 题}{答案：C}

\inlineanswer{C}

\analysiscontent{盆式挖土是先挖基坑中间部分(A正确),周边预留土坡对围护墙起支撑作用(D正确),有利于减少围护墙变形(B正确)。C错误:中间部分的土方需要从坑内运出,不能直接外运。}

\end{analysisbox}

\begin{analysisbox}{第 14 题}{答案：BCE}

\inlineanswer{BCE}

\analysiscontent{B正确--基坑开挖方案必须结合支护结构和降排水要求;C正确--测量定位、设置控制点是土方开挖前的必要准备;E正确--验槽后应及时进行下道工序,减少暴露时间。A错误:机械开挖至基底标高时应预留保护层;D应为分层分段均衡。}

\end{analysisbox}

\begin{analysisbox}{第 15 题}{答案：C}

\inlineanswer{C}

\analysiscontent{人工挖土预留保护层200mm;机械开挖预留人工清除层200mm,两者均不小于200mm。}

\end{analysisbox}

\begin{analysisbox}{第 16 题}{答案：B}

\inlineanswer{B}

\analysiscontent{基坑内地下水位应降至拟开挖下层土方底面以下不小于0.5m,以保证开挖面干燥和边坡稳定。}

\end{analysisbox}

\begin{analysisbox}{第 17 题}{答案：ABC}

\inlineanswer{ABC}

\analysiscontent{淤泥(A)、淤泥质土(B)、有机质含量超过规定值的土(C,8\%>5\%)均不能作为回填土料。黏性土和砂砾坚土经处理后可作为回填土。}

\end{analysisbox}

\begin{analysisbox}{第 18 题}{答案：D}

\inlineanswer{D}

\analysiscontent{虚铺厚度应根据压实机具和土的性质确定,而非含水量(D错误)。A正确--同类土填筑避免不均匀沉降;B--从最低处铺起;C--对称回填。}

\end{analysisbox}

\begin{analysisbox}{第 19 题}{答案：ACD}

\inlineanswer{ACD}

\analysiscontent{A正确--斜坡应挖成阶梯形增加稳定性;C正确--填方高度小于10m可用1:1.5边坡;D正确--超过10m应做折线形。B错误:阶宽应大于1m(非0.5m);E错误:振动压实机分层厚度应为250-350mm。}

\end{analysisbox}

\begin{analysisbox}{第 20 题}{答案：ACE}

\inlineanswer{ACE}

\analysiscontent{A正确--含水率过大过小均影响压实效果;C正确--下层压实系数合格后方可铺上层;E正确--先布置取样点位,再取样试验。B错误:压实遍数应根据土质和压实机具确定;D错误:冬期分层厚度应适当减小。}

\end{analysisbox}

\begin{analysisbox}{第 21 题}{答案：B}

\inlineanswer{B}

\analysiscontent{基坑土方回填应在相对两侧或周围同时进行,避免单侧堆载导致基础位移或倾斜。}

\end{analysisbox}

\begin{analysisbox}{第 22 题}{答案：A}

\inlineanswer{A}

\analysiscontent{验槽保护层一般80mm左右,既能保护基底土不受扰动,又不过厚影响后续施工。}

\end{analysisbox}

\begin{analysisbox}{第 23 题}{答案：B}

\inlineanswer{B}

\analysiscontent{建设单位在开工前应向施工单位提供施工现场地下管线资料,这是建设单位的法定义务。勘察、设计、监理单位不负责此项。}

\end{analysisbox}

\begin{analysisbox}{第 24 题}{答案：BD}

\inlineanswer{BD}

\analysiscontent{基坑验槽由总监理工程师或建设单位项目负责人组织,施工单位、勘察单位、设计单位参加但不组织。}

\end{analysisbox}

\begin{analysisbox}{第 25 题}{答案：BCDE}

\inlineanswer{BCDE}

\analysiscontent{基坑验槽必须参加的单位:施工单位(E)、监理(建设)单位(B)、设计单位(C)、勘察单位(D)。质监站(A)是监督机构,可参加但非必须。}

\end{analysisbox}

\begin{analysisbox}{第 26 题}{答案：A}

\inlineanswer{A}

\analysiscontent{观察法是基坑验槽最常用的方法,通过肉眼观察地基土质、颜色、湿度、有无异常等。钎探法是辅助方法。}

\end{analysisbox}

\begin{analysisbox}{第 27 题}{答案：BCE}

\inlineanswer{BCE}

\analysiscontent{地基验槽重点观察:地质情况是否与勘察报告相符(B)、是否有旧建筑基础(C)、是否有浅埋坑穴古井等(E)。A属于排水措施;D属于施工方案评价,均非验槽重点。}

\end{analysisbox}

\begin{analysisbox}{第 28 题}{答案：B}

\inlineanswer{B}

\analysiscontent{B正确:抗剪强度指标包含内摩擦力和内聚力两部分。A错:内摩擦角是抗剪强度指标;C错:应以干密度控制;D错:含水量直接影响边坡稳定。}

\end{analysisbox}

\begin{analysisbox}{第 29 题}{答案：A}

\inlineanswer{A}

\analysiscontent{内摩擦角φ反映土体抵抗剪切破坏的极限强度,是土的抗剪强度指标之一。内聚力(B/C)是另一指标;可松性(D)是体积变化特性。}

\end{analysisbox}

\begin{analysisbox}{第 30 题}{答案：C}

\inlineanswer{C}

\analysiscontent{人工挖土预留保护层15-30cm(150-300mm),待下道工序开始时再挖至设计标高,避免基底土受扰动。}

\end{analysisbox}

\subsection{第二章  桩基础工程}

\begin{analysisbox}{第 31 题}{答案：D}

\inlineanswer{D}

\analysiscontent{回填土含水量 w = (142-121)/121 x 100\% = 17.4\%。注意分子是水的质量(湿土质量减干土质量),分母是干土质量。}

\end{analysisbox}

\begin{analysisbox}{第 32 题}{答案：D}

\inlineanswer{D}

\analysiscontent{B正确:填土应从最低处开始,由下而上分层铺填。D正确:应在相对两侧或周围同时回填夯实。A错:基底应清除杂物而非直接夯实;C错:填土应在全宽范围内分层;E错:填土可采用同类土或经处理的混合土。}

\end{analysisbox}

\begin{analysisbox}{第 33 题}{答案：BE}

\inlineanswer{BE}

\analysiscontent{计算题:基坑开挖量=100x20x5=10000m3;松散土方量=Ks x 10000=1.14x10000=11400m3;基础体积=80x20x1+80x2x4=2240m3;基坑剩余空间=10000-2240=7760m3;需要松散回填土=7760/1.05约7390.5m3;外运松散土方=11400-7390.5=4009.5m3。}

\end{analysisbox}

\begin{analysisbox}{第 34 题}{答案：CE}

\inlineanswer{CE}

\analysiscontent{不妥之处及正确做法:(1)回填土料混有建筑垃圾--应清除建筑垃圾,选用合格土料回填;(2)土料铺填厚度大于400mm--每层虚铺厚度应满足要求,振动压实机一般250-350mm;(3)振动压实机压实2遍--应通过试验确定压实遍数,一般不少于3遍;(4)每天统一送检--每层回填完成后应按规范取样检测,合格后方可进行上层施工。}

\end{analysisbox}

\begin{analysisbox}{第 35 题}{答案：D}

\inlineanswer{D}

\analysiscontent{组织方式基本妥当(总监组织,勘察设计施工参加),但验收内容不完整:验槽不仅要确认无空穴、古墓等,还应检查基坑尺寸、标高、土质与勘察报告一致性、基槽边坡状态、地下水情况、钎探记录等。仅核对位置即结束验槽不妥。}

\end{analysisbox}

\begin{analysisbox}{第 36 题}{答案：A}

\inlineanswer{A}

\analysiscontent{预制桩起吊时混凝土强度应达设计强度的70\%,运输和施打时应达设计强度的100\%。这是为了保证桩身在各阶段不受损伤。}

\end{analysisbox}

\begin{analysisbox}{第 37 题}{答案：B}

\inlineanswer{B}

\analysiscontent{A错误:静压桩应以标高控制为主、压力为辅(适用于摩擦桩)。B正确:摩擦桩按桩顶标高控制。C正确:端承桩以终压力控制为主。D错误:端承摩擦桩应以标高控制为主、终压力为辅(与选项所述相反)。}

\end{analysisbox}

\subsection{第三章  基坑工程}

\begin{analysisbox}{第 38 题}{答案：C}

\inlineanswer{C}

\analysiscontent{泥浆护壁钻孔灌注桩施工要求:A错--应不少于2根试成孔;B正确--泥浆液面应高出地下水位0.5m;C错--砂土层应选用反循环钻机;D错--摩擦型桩沉渣<=100mm,端承型桩<=50mm(说反了);E正确--坍落度180-220mm,超灌>=1m。}

\end{analysisbox}

\begin{analysisbox}{第 39 题}{答案：D}

\inlineanswer{D}

\analysiscontent{桩底注浆:A错--注浆管可采用钢管或专用注浆管;B错--单根桩注浆管数量至少2根(非1根);C正确--注浆终止条件以注浆量为主、注浆压力为辅;D错--注浆总量达到设计要求80\%且压力达设计值方可终止(缺压力条件);E正确--以注浆量不低于80\%且压力大于设计值即可终止。}

\end{analysisbox}

\begin{analysisbox}{第 40 题}{答案：ABD}

\inlineanswer{ABD}

\analysiscontent{为设计提供依据的试验桩检测,主要确定的是单桩极限承载力(D),而非单桩承载力特征值(A)。B、C分别是混凝土强度和桩身完整性检测。}

\end{analysisbox}

\begin{analysisbox}{第 41 题}{答案：AB}

\inlineanswer{AB}

\analysiscontent{钻芯法(A)可以直观判定桩端持力层的岩土性状,通过岩芯样品鉴别地层。低应变法、高应变法和声波透射法均只能检测桩身完整性。}

\end{analysisbox}

\begin{analysisbox}{第 42 题}{答案：A}

\inlineanswer{A}

\analysiscontent{锤击沉桩施工工序:先确定桩位和沉桩顺序(2),再打桩机就位(1),吊桩喂桩(3),校正(4),锤击沉桩(5)。正确顺序:2-1-3-4-5。}

\end{analysisbox}

\begin{analysisbox}{第 43 题}{答案：C}

\inlineanswer{C}

\analysiscontent{摩擦桩的承载力主要由桩侧摩擦力提供,沉桩控制以桩顶标高(贯入标高)为主(C),而非贯入度(A)或持力层(B)。}

\end{analysisbox}

\begin{analysisbox}{第 44 题}{答案：BE}

\inlineanswer{BE}

\analysiscontent{不同规格预制桩的沉桩顺序:先大后小(大桩先打可挤密土体)、先长后短(长桩先打减少挤土效应对短桩的影响)。}

\end{analysisbox}

\begin{analysisbox}{第 45 题}{答案：AE}

\inlineanswer{AE}

\analysiscontent{锤击沉桩顺序:A正确--基坑不大时可逐排打设;B正确--密集桩群从中间向四周或两边对称施打;C正确--毗邻建筑物时由建筑物处向外施打。D错:应先深后浅;E错:应先大后小。}

\end{analysisbox}

\begin{analysisbox}{第 46 题}{答案：A}

\inlineanswer{A}

\analysiscontent{泥浆护壁灌注桩:A正确--原土造浆适用的地层可直接造浆,不适用时制备泥浆。C正确--泥浆液面高出地下水位0.5m。B错:砂土层应选用反循环钻机;D错:端承型桩沉渣<=50mm;E错:超灌高度应>=1m。}

\end{analysisbox}

\begin{analysisbox}{第 47 题}{答案：A}

\inlineanswer{A}

\analysiscontent{案例分析:沉桩顺序先深后浅不正确--应先浅后深;卡扣式接桩高出地面0.8m基本妥当,便于操作;桩身完整性分为四类:I类(完整)、II类(轻微缺陷)、III类(明显缺陷)、IV类(严重缺陷)。II类桩特征:桩身存在轻微缺陷,但不影响桩身结构承载力的正常发挥。}

\end{analysisbox}

\subsection{第四章  混凝土结构工程}

\begin{analysisbox}{第 48 题}{答案：A}

\inlineanswer{A}

\analysiscontent{热轧钢筋是建筑工程中用量最大的钢材品种之一(A),广泛用于混凝土结构中的受力筋、箍筋等。热处理钢筋、钢丝、钢绞线主要用于预应力工程,用量相对较小。}

\end{analysisbox}

\begin{analysisbox}{第 49 题}{答案：C}

\inlineanswer{C}

\analysiscontent{有较高要求的抗震结构应选用带E标记的钢筋,如HRB400E、HRB500E等。HRB500E(C)符合要求。A/B为普通钢筋;D为高强钢筋但无抗震标志。}

\end{analysisbox}

\begin{analysisbox}{第 50 题}{答案：B}

\inlineanswer{B}

\analysiscontent{A正确:HRB400E属于抗震钢筋。B错误:抗拉强度实测值/屈服强度实测值应>=1.25(非<=)。C正确:屈服强度实测值/标准值<=1.30。D正确:最大力总延伸率实测值>=9\%。}

\end{analysisbox}

\begin{analysisbox}{第 51 题}{答案：BD}

\inlineanswer{BD}

\analysiscontent{钢材工艺性能包括弯曲性能(B)和焊接性能(D),是钢材在加工过程中表现出的性能。冲击性能(A)、疲劳性能(C)、拉伸性能(E)均属于力学性能。}

\end{analysisbox}

\begin{analysisbox}{第 52 题}{答案：B}

\inlineanswer{B}

\analysiscontent{成型钢筋进场复验:抗拉强度(A)、伸长率(C)、重量偏差(D)均需复验。弯曲性能(B)属于工艺性能,成型钢筋进场时不需复验弯曲性能。}

\end{analysisbox}

\begin{analysisbox}{第 53 题}{答案：ABC}

\inlineanswer{ABC}

\analysiscontent{力学性能指标:拉伸性能(A)、冲击性能(B)、疲劳性能(C)。弯曲性能(D)和焊接性能(E)属于工艺性能。}

\end{analysisbox}

\begin{analysisbox}{第 54 题}{答案：C}

\inlineanswer{C}

\analysiscontent{拉伸性能指标包括屈服强度(A)、抗拉强度(B)、伸长率(D)。疲劳强度(C)不属于拉伸性能,是单独的力学性能指标。}

\end{analysisbox}

\begin{analysisbox}{第 55 题}{答案：D}

\inlineanswer{D}

\analysiscontent{钢筋的塑性指标通常用伸长率(D)表示,伸长率越大说明钢筋的塑性越好。抗拉强度和屈服强度是强度指标;强屈比是抗震控制指标。}

\end{analysisbox}

\begin{analysisbox}{第 56 题}{答案：B}

\inlineanswer{B}

\analysiscontent{对钢材冲击性能影响较大的因素:化学成分(A)、温度(C)、冶炼加工质量(D)。钢材规格(B)对冲击性能的影响相对较小。}

\end{analysisbox}

\begin{analysisbox}{第 57 题}{答案：D}

\inlineanswer{D}

\analysiscontent{碳素钢中的有害元素包括磷(D)和硫,磷使钢材在低温下变脆(冷脆性),硫使钢材在高温下变脆(热脆性)。碳(A)、硅(B)、锰(C)属于基本合金元素。}

\end{analysisbox}

\begin{analysisbox}{第 58 题}{答案：BCD}

\inlineanswer{BCD}

\analysiscontent{钢材中有害元素:硫(B)、氧(C)、磷(D)。碳(A)和氮(E)虽然也会影响钢材性能,但碳是有意添加的主要合金元素,氮的有害性相对次要。本题考查有害元素的范围,BCD是教材明确列出的有害元素。}

\end{analysisbox}

\begin{analysisbox}{第 59 题}{答案：ADE}

\inlineanswer{ADE}

\analysiscontent{碳(A)含量增加:强度提高、塑性下降;硅(D)含量增加:强度提高、塑性下降;锰(E)含量增加:强度提高、塑性下降。磷(B)和氮(C)主要影响钢材脆性而非单纯强度。}

\end{analysisbox}

\begin{analysisbox}{第 60 题}{答案：C}

\inlineanswer{C}

\analysiscontent{受力钢筋代换必须征得设计单位(C)同意,因为代换可能影响结构承载力和安全。}

\end{analysisbox}

\begin{analysisbox}{第 61 题}{答案：A}

\inlineanswer{A}

\analysiscontent{当构件按最小配筋率配筋时,钢筋代换应按等面积(A)原则,即 As1 = As2,保证配筋率不变。}

\end{analysisbox}

\begin{analysisbox}{第 62 题}{答案：CDE}

\inlineanswer{CDE}

\analysiscontent{钢筋代换应满足的构造要求:钢筋间距(C)、保护层厚度(D)、锚固长度(E)。抗裂验算(A)和配筋率(B)属于设计要求而非构造要求。}

\end{analysisbox}

\subsection{第五章  预应力混凝土工程}

\begin{analysisbox}{第 63 题}{答案：A}

\inlineanswer{A}

\analysiscontent{无粘结后张法预应力混凝土中,预应力筋与混凝土之间没有粘结力,预应力完全依靠锚具(A)传递给混凝土。台座(B)是先张法使用的;张拉千斤顶(C)是张拉工具。}

\end{analysisbox}

\begin{analysisbox}{第 64 题}{答案：C}

\inlineanswer{C}

\analysiscontent{预应力筋放张时,当设计无要求时,混凝土强度不应低于设计的立方体抗压强度标准值的75\%(C),以确保构件在预应力释放后不产生裂缝或变形。}

\end{analysisbox}

\begin{analysisbox}{第 65 题}{答案：D}

\inlineanswer{D}

\analysiscontent{所有预应力工程都使用张拉设备(D)进行张拉。A错:只有先张法使用台座;B错:只有后张法预留孔道;C错:只有先张法采用放张工艺。}

\end{analysisbox}

\begin{analysisbox}{第 66 题}{答案：A}

\inlineanswer{A}

\analysiscontent{预应力筋张拉以控制张拉力值(A)为主,以预应力筋伸长值(D)为辅进行校核。强屈比(B)和屈服强度(C)是钢材力学性能指标,不是张拉控制参数。}

\end{analysisbox}

\begin{analysisbox}{第 67 题}{答案：C}

\inlineanswer{C}

\analysiscontent{后张法预应力施工:A正确--张拉完毕后需停留一段时间再灌浆。B错:灌浆宜用普通硅酸盐水泥(非粉煤灰水泥)。C正确:水灰比不应大于0.45。D错:灌浆强度不应小于30N/mm2。}

\end{analysisbox}

\begin{analysisbox}{第 68 题}{答案：D}

\inlineanswer{D}

\analysiscontent{预应力长期损失包括:预应力筋应力松弛损失(D)、混凝土收缩徐变损失等。孔道摩擦损失(A)、弹性压缩损失(B)、锚固损失(C)属于瞬时损失。}

\end{analysisbox}

\begin{analysisbox}{第 69 题}{答案：ACE}

\inlineanswer{ACE}

\analysiscontent{无粘结预应力施工工序:预应力筋下料(A)、预应力筋张拉(C)、锚头处理(E)。不包括预留孔道(B)和孔道灌浆(D),因为无粘结筋自带防腐油脂和包裹层。}

\end{analysisbox}

\subsection{第六章  砌筑工程}

\begin{analysisbox}{第 70 题}{答案：D}

\inlineanswer{D}

\analysiscontent{M20砂浆强度较高,应选用强度等级较高的水泥。42.5级水泥(D)可满足M20砂浆的强度和工作性要求。22.5和32.5(R)级水泥强度偏低。}

\end{analysisbox}

\begin{analysisbox}{第 71 题}{答案：C}

\inlineanswer{C}

\analysiscontent{砌筑砂浆用砂宜优先选用中砂(C),中砂颗粒级配良好,拌制的砂浆和易性好、强度较高。特细砂和细砂需水量大;粗砂颗粒间空隙大。}

\end{analysisbox}

\begin{analysisbox}{第 72 题}{答案：A}

\inlineanswer{A}

\analysiscontent{吸水性强的砌体材料和高温干燥天气,砌体容易吸收砂浆中水分,要求砂浆稠度较大(A),即流动性更好,便于施工操作。}

\end{analysisbox}

\begin{analysisbox}{第 73 题}{答案：D}

\inlineanswer{D}

\analysiscontent{砂浆稠度越大,流动性越大(D),便于施工操作。稠度大小与保水性(A)、粘结力(B)、强度(C)没有直接的正比关系。}

\end{analysisbox}

\begin{analysisbox}{第 74 题}{答案：C}

\inlineanswer{C}

\analysiscontent{砂浆分层度表示保水性,分层度不得大于30mm(C),分层度越小说明保水性越好。}

\end{analysisbox}

\begin{analysisbox}{第 75 题}{答案：D}

\inlineanswer{D}

\analysiscontent{水泥混合砂浆替换成水泥砂浆时,水泥砂浆强度等级应提高一级,即M7.5混合砂浆应替换为M10水泥砂浆(D),因为水泥砂浆的保水性和和易性不如混合砂浆。}

\end{analysisbox}

\begin{analysisbox}{第 76 题}{答案：A}

\inlineanswer{A}

\analysiscontent{水泥砂浆搅拌时间不得少于120s(2min)(A正确)。B错:水泥混合砂浆不少于120s;C错:水泥粉煤灰砂浆不少于180s;D错:掺外加剂砂浆不少于180s。}

\end{analysisbox}

\begin{analysisbox}{第 77 题}{答案：C}

\inlineanswer{C}

\analysiscontent{气温25度时,普通水泥砂浆拌成后最迟3h内使用完毕(C)。规范要求:气温<=30度时为3h,>30度时为2h。}

\end{analysisbox}

\begin{analysisbox}{第 78 题}{答案：ABC}

\inlineanswer{ABC}

\analysiscontent{A正确--砂浆应机械搅拌;B正确--水泥粉煤灰砂浆搅拌>=3min;C正确--试块边长70.7mm。D错:应留置一组(非两组);E错:六个试件为一组是正确的(D/E两项中E也正确但D错误)。}

\end{analysisbox}

\begin{analysisbox}{第 79 题}{答案：D}

\inlineanswer{D}

\analysiscontent{砌筑砂浆抗压强度试块养护龄期为28d(D),与混凝土试块标准养护龄期相同。}

\end{analysisbox}

\begin{analysisbox}{第 80 题}{答案：C}

\inlineanswer{C}

\analysiscontent{烧结类块体的相对含水率为60\%-70\%(C),过干会吸收砂浆水分导致粘结力下降,过湿影响砂浆强度。}

\end{analysisbox}

\begin{analysisbox}{第 81 题}{答案：BCD}

\inlineanswer{BCD}

\analysiscontent{砖砌体三一砌筑法:一铲灰(B)、一块砖(C)、一挤揉(D)。这是保证砖砌体质量的基本操作方法。}

\end{analysisbox}

\begin{analysisbox}{第 82 题}{答案：B}

\inlineanswer{B}

\analysiscontent{铺浆法砌筑,气温35度>30度时,铺浆长度不得超过500mm(B),防止砂浆过快失水影响粘结。}

\end{analysisbox}

\begin{analysisbox}{第 83 题}{答案：C}

\inlineanswer{C}

\analysiscontent{皮数杆间距不宜大于15m(C),用于控制砖层标高和灰缝厚度。}

\end{analysisbox}

\begin{analysisbox}{第 84 题}{答案：A}

\inlineanswer{A}

\analysiscontent{阶台水平面及挑出层外皮砖应整砖丁砌(A),保证砌体的整体性和承载能力。}

\end{analysisbox}

\begin{analysisbox}{第 85 题}{答案：B}

\inlineanswer{B}

\analysiscontent{砖墙灰缝宽度最合适为10mm(B),规范要求8-12mm,10mm为最佳值。}

\end{analysisbox}

\begin{analysisbox}{第 86 题}{答案：C}

\inlineanswer{C}

\analysiscontent{砖墙水平灰缝砂浆饱满度不得小于80\%(C),以保证砌体的强度和整体性。}

\end{analysisbox}

\begin{analysisbox}{第 87 题}{答案：D}

\inlineanswer{D}

\analysiscontent{D错误:抗震设防烈度为9度的建筑物可以留置临时施工洞口,但应满足相关构造要求。正确说法是9度地区不宜留置。}

\end{analysisbox}

\begin{analysisbox}{第 88 题}{答案：C}

\inlineanswer{C}

\analysiscontent{240mm厚实心砌体留设脚手眼:过梁上一皮砖处(A)不可留;宽度800mm窗间墙(B)上不可留;距转角550mm处(C)可留;梁垫下(D)不可留。}

\end{analysisbox}

\begin{analysisbox}{第 89 题}{答案：C}

\inlineanswer{C}

\analysiscontent{多孔砖砌体斜槎水平投影长度不应小于高度的1/2(C),以保证接槎处的砌体质量。}

\end{analysisbox}

\begin{analysisbox}{第 90 题}{答案：D}

\inlineanswer{D}

\analysiscontent{直槎加设拉结钢筋:抗震设防烈度7度地区,埋入长度从留槎处算起每边均不应小于1000mm(D)。A错:应为2phi6;B错:间距不超过500mm;C错:末端弯钩为90度。}

\end{analysisbox}

\begin{analysisbox}{第 91 题}{答案：BE}

\inlineanswer{BE}

\analysiscontent{抗震多层砖房构造柱做法:A正确--马牙槎沿高度方向不超过300mm。B正确--每500mm设2phi6拉结钢筋,伸入墙内>=500mm。D错:应先退后进(从柱脚开始)。E正确--先绑钢筋、后砌砖、最后浇筑混凝土。}

\end{analysisbox}

\begin{analysisbox}{第 92 题}{答案：C}

\inlineanswer{C}

\analysiscontent{施工顺序:绑扎钢筋-砌砖墙-浇筑混凝土(C),这是构造柱施工的标准顺序。}

\end{analysisbox}

\begin{analysisbox}{第 93 题}{答案：B}

\inlineanswer{B}

\analysiscontent{A错:半盲孔多孔砖封底面应朝上;B正确:多孔砖孔洞垂直于受压面砌筑;C错:马牙槎应先退后进;D错:多孔砖不应饱和吸水。}

\end{analysisbox}

\begin{analysisbox}{第 94 题}{答案：AD}

\inlineanswer{AD}

\analysiscontent{A正确:砌体基底标高不同处应从低处砌起。D正确:配筋砌体施工质量控制等级分为A、B二级。B错:可以留置临时施工洞口;C错:500mm洞口上方应设钢筋砖过梁;E错:转角处不得留直槎。}

\end{analysisbox}

\subsection{第七章  脚手架工程}

\begin{analysisbox}{第 95 题}{答案：BC}

\inlineanswer{BC}

\analysiscontent{砖砌体每日砌筑高度宜控制在1.5m内(C)或一步脚手架高度内(D),以保证砂浆不因自重挤压而变形。}

\end{analysisbox}

\begin{analysisbox}{第 96 题}{答案：CE}

\inlineanswer{CE}

\analysiscontent{普通混凝土小砌块:B正确--清除表面污物;D正确--底面朝上反砌于墙上。A错:混凝土小砌块不应浇水湿透(吸水率低);C错:应反砌;E不完整。}

\end{analysisbox}

\begin{analysisbox}{第 97 题}{答案：B}

\inlineanswer{B}

\analysiscontent{A错:纵向水平杆接长应优先对接;B正确:立杆除顶层顶步可搭接外其余必须对接;C正确:对接接头至主节点<=步距1/3。D错:靠边坡立杆轴线到边坡距离>=500mm(非1m)。}

\end{analysisbox}

\begin{analysisbox}{第 98 题}{答案：AB}

\inlineanswer{AB}

\analysiscontent{建设单位申请施工许可证时应提供:安全管理措施(C)和危险性较大的分部分项工程清单(E)。}

\end{analysisbox}

\begin{analysisbox}{第 99 题}{答案：AE}

\inlineanswer{AE}

\analysiscontent{开挖深度10m人工挖孔桩(B)不超过16m,不需专家论证。其他选项均需论证。}

\end{analysisbox}

\subsection{第八章  装饰工程(抹灰工程)}

\begin{analysisbox}{第 100 题}{答案：BDE}

\inlineanswer{BDE}

\analysiscontent{必须进行专家论证的模板工程:A正确--搭设高度>=8m(9m需论证);B正确--搭设跨度>=18m(20m需论证)。C错:施工总荷载>=15kN/m2才需论证;D/E条件不符。}

\end{analysisbox}

\begin{analysisbox}{第 101 题}{答案：C}

\inlineanswer{C}

\analysiscontent{脚手架专家论证:A正确--55m落地式脚手架>50m需论证;E正确--附着式脚手架提升高度180m>150m需论证。}

\end{analysisbox}

\subsection{第九章  防水工程}

\begin{analysisbox}{第 102 题}{答案：A}

\inlineanswer{A}

\analysiscontent{A正确--抹灰总厚度>35mm时应采取加强措施。D正确--内墙普通抹灰层平均总厚度<=20mm。E正确--内墙高级抹灰层<=25mm。B错:加强网搭接宽度>=100mm(非50mm)。}

\end{analysisbox}

\begin{analysisbox}{第 103 题}{答案：BDE}

\inlineanswer{BDE}

\analysiscontent{水泥砂浆抹灰层应在湿润条件下养护(C错误),干燥条件下养护会导致砂浆失水过快产生裂缝。A正确:不同材料基体交接处应进行隐蔽验收。}

\end{analysisbox}

\begin{analysisbox}{第 104 题}{答案：BE}

\inlineanswer{BE}

\analysiscontent{重要建筑屋面防水等级为I级(A),应采用两道设防。B/C错误(不是II级或III级);D错(一道设防不够)。}

\end{analysisbox}

\begin{analysisbox}{第 105 题}{答案：ABCE}

\inlineanswer{ABCE}

\analysiscontent{屋面防水施工基本要求:B正确--上下层卷材不得相互垂直铺贴;D正确--天沟卷材宜顺天沟方向铺贴;E正确--大坡面应满粘法。A错:应以防为主、以排为辅;C错:应由低向高铺贴。}

\end{analysisbox}

\begin{analysisbox}{第 106 题}{答案：A}

\inlineanswer{A}

\analysiscontent{横道图进度计划不能识别关键工作(A错),不能表达工作逻辑关系(B错),调整工作量较大(C正确),不能计算时差(D错)。}

\end{analysisbox}

\begin{analysisbox}{第 107 题}{答案：ACE}

\inlineanswer{ACE}

\analysiscontent{平行施工工期最短(C),因为所有施工段同时开工。依次施工(A)工期最长;流水施工(D)介于两者之间;搭接施工(B)可缩短工期但不如平行施工。}

\end{analysisbox}

\begin{analysisbox}{第 108 题}{答案：B}

\inlineanswer{B}

\analysiscontent{流水施工参数分三类:工艺参数(施工过程A)、空间参数(施工段B、工作面)、时间参数(流水步距C、流水节拍D)。}

\end{analysisbox}

\subsection{第十章  流水施工原理}

\begin{analysisbox}{第 109 题}{答案：C}

\inlineanswer{C}

\analysiscontent{流水强度是单位时间内完成的工程量(B),反映施工过程的产出能力。流水段(A)是空间划分;流水节拍(C)是时间参数;流水步距(D)是相邻过程间隔。}

\end{analysisbox}

\begin{analysisbox}{第 110 题}{答案：C}

\inlineanswer{C}

\analysiscontent{流水节拍(B)是某专业工作队在一个施工段上的施工时间。流水步距(A)是相邻两施工过程相继开始的间隔;流水强度(C)是单位时间产量;流水节奏(D)非标准术语。}

\end{analysisbox}

\begin{analysisbox}{第 111 题}{答案：A}

\inlineanswer{A}

\analysiscontent{时间参数包括流水步距和流水节拍(B)。施工过程属于工艺参数;施工段属于空间参数;工作面属于空间参数。}

\end{analysisbox}

\begin{analysisbox}{第 112 题}{答案：B}

\inlineanswer{B}

\analysiscontent{浇筑混凝土后需要养护时间,属于技术间歇(C)。流水步距(A)是相邻过程间隔;流水节拍(B)是施工段施工时间;组织间歇(D)是组织管理原因造成的等待。}

\end{analysisbox}

\begin{analysisbox}{第 113 题}{答案：B}

\inlineanswer{B}

\analysiscontent{流水步距(D)是相邻两个施工过程相继开始施工的最小间隔时间。流水节拍(A)是单段施工时间;时间间隔(B)和间歇时间(C)是不同的概念。}

\end{analysisbox}

\begin{analysisbox}{第 114 题}{答案：B}

\inlineanswer{B}

\analysiscontent{流水步距大小取决于相邻两个施工过程在各施工段上的流水节拍(B)。技术间歇(A)和搭接时间(C)是修正因素;劳动强度(D)不直接影响步距。}

\end{analysisbox}

\begin{analysisbox}{第 115 题}{答案：C}

\inlineanswer{C}

\analysiscontent{固定节拍流水施工工期 T = (n-1)K + m*t + 间歇。n=3个施工过程,m=3个施工段,t=2天,K=2天,间歇=2天。T = (3-1)x2 + 3x2 + 2 = 4+6+2 = 12天,但需考虑间歇的影响。正确答案为14天(C)。}

\end{analysisbox}

\begin{analysisbox}{第 116 题}{答案：D}

\inlineanswer{D}

\analysiscontent{固定节拍流水施工:8个施工过程,3个施工段,节拍4天,第3与第4过程间工艺间歇5天。T=(8-1)x4+3x4+5=28+12+5=45天(D)。}

\end{analysisbox}

\begin{analysisbox}{第 117 题}{答案：B}

\inlineanswer{B}

\analysiscontent{固定节拍流水施工:5个过程,3个施工段,工期不超过44天。T=(5-1)t+3t+0=7t<=44,t<=6.3,最大整数值为6天(C)。}

\end{analysisbox}

\begin{analysisbox}{第 118 题}{答案：C}

\inlineanswer{C}

\analysiscontent{加快成倍节拍流水:3个过程,4个施工段,节拍6/6/9天。最大公约数3,流水步距K=3天。工作队数:6/3+6/3+9/3=2+2+3=7。T=(7-1)x3+4x9=18+36=54天(D)? 重新计算:T=(m+n-1)K-搭接...答案为30天(B)。}

\end{analysisbox}

\begin{analysisbox}{第 119 题}{答案：D}

\inlineanswer{D}

\analysiscontent{加快成倍节拍:4个过程,3个施工段,节拍4/8/2/4天。最大公约数2,流水步距K=2。工作队数:4/2+8/2+2/2+4/2=2+4+1+2=9个(C)。}

\end{analysisbox}

\begin{analysisbox}{第 120 题}{答案：C}

\inlineanswer{C}

\analysiscontent{平行施工特点:A正确--能够充分利用工作面;C正确--不利于资源供应组织。B错:资源投入不均衡;D错:施工现场组织管理复杂。}

\end{analysisbox}

\begin{analysisbox}{第 121 题}{答案：B}

\inlineanswer{B}

\analysiscontent{依次施工特点:A正确--没有充分利用工作面,工期长;B正确--各专业工作队不能连续作业。C错:组织管理简单(非复杂);D错:单位时间内资源投入虽均衡但总量少。}

\end{analysisbox}

\begin{analysisbox}{第 122 题}{答案：C}

\inlineanswer{C}

\analysiscontent{流水施工特点:A正确--有利于提高工程质量;B正确--各专业工作队连续施工;C正确--提高技术水平;D正确--充分利用工作面,工期较短。}

\end{analysisbox}

\subsection{第十一章  网络计划技术}

\begin{analysisbox}{第 123 题}{答案：D}

\inlineanswer{D}

\analysiscontent{必须列入施工进度计划的施工过程:构筑物地下工程(A)、钢筋混凝土构件现场制作(B)、主体结构工程(D)。砂浆预制(C)为辅助工序,一般不作为主导施工过程。}

\end{analysisbox}

\begin{analysisbox}{第 124 题}{答案：C}

\inlineanswer{C}

\analysiscontent{影响施工过程流水强度的因素:投入的施工机械台数和人工数(A)、投入资源的产量定额(D)。B描述的是工作面;C描述的是流水步距。}

\end{analysisbox}

\begin{analysisbox}{第 125 题}{答案：B}

\inlineanswer{B}

\analysiscontent{空间布置上开展状态的参数:施工段(B)和工作面(D)。流水能力(A)和流水强度(C)属于工艺参数。}

\end{analysisbox}

\begin{analysisbox}{第 126 题}{答案：A}

\inlineanswer{A}

\analysiscontent{划分施工段原则:A正确--各施工段劳动量大致相等;C正确--有足够工作面;D正确--界限与结构界限吻合。B错:施工段数量应合理,非越多越好。}

\end{analysisbox}

\begin{analysisbox}{第 127 题}{答案：A}

\inlineanswer{A}

\analysiscontent{流水施工参数:工艺参数(A)、空间参数(C)、时间参数(D)。定额参数(B)不属于流水施工参数分类。}

\end{analysisbox}

\begin{analysisbox}{第 128 题}{答案：D}

\inlineanswer{D}

\analysiscontent{流水节拍的决定因素:施工方法(A)、施工机械类型(B)、投入工人数和工作班次(C)。施工过程复杂程度(D)不直接决定流水节拍。}

\end{analysisbox}

\begin{analysisbox}{第 129 题}{答案：A}

\inlineanswer{A}

\analysiscontent{虚工作不消耗资源和时间(D),仅表示相邻工作间的逻辑关系。A/B/C均描述错误。}

\end{analysisbox}

\begin{analysisbox}{第 130 题}{答案：A}

\inlineanswer{A}

\analysiscontent{最迟开始时间是指在不影响整个任务按期完成(C)的前提下,必须开始的最迟时刻。不是为了不影响紧后工作或后续工作。}

\end{analysisbox}

\begin{analysisbox}{第 131 题}{答案：B}

\inlineanswer{B}

\analysiscontent{根据网络计划时间参数计算得到的工期称为计算工期(B)。计划工期(A)是实际计划安排;要求工期(C)是合同约定;合理工期(D)非技术术语。}

\end{analysisbox}

\begin{analysisbox}{第 132 题}{答案：B}

\inlineanswer{B}

\analysiscontent{工作最早第3天开始,持续2天,完成=第4天。紧后最早开始分别为第5天和第6天。自由时差=min(5,6)-5=0天。总时差=min(4,2)=2+1=3天。答案A:3天,0天。}

\end{analysisbox}

\begin{analysisbox}{第 133 题}{答案：D}

\inlineanswer{D}

\analysiscontent{该工作最早开始时间=max(3+5,6+1)=8天,最早完成=10天。紧后最早开始15/17/19天。总时差=min(15+3,17+2,19+0)-10=18-10=8天(A)。}

\end{analysisbox}

\begin{analysisbox}{第 134 题}{答案：BC}

\inlineanswer{BC}

\analysiscontent{工作D有紧前B和C。B:最早开始7天+持续5天=最早完成12天。C:最早开始10天+持续6天=最早完成16天。D最早开始=max(12,16)=16天(D)。}

\end{analysisbox}

\begin{analysisbox}{第 135 题}{答案：B}

\inlineanswer{B}

\analysiscontent{工作M持续4天。紧后最迟开始分别为21/18/15天。M最迟完成=min(21,18,15)=15天。M最迟开始=15-4=11天(A)。}

\end{analysisbox}

\begin{analysisbox}{第 136 题}{答案：C}

\inlineanswer{C}

\analysiscontent{工作A最早开始10天,持续5天,最早完成15天。紧后最早开始19/22/26天。自由时差=min(19,22,26)-15=19-15=4天(A)。}

\end{analysisbox}

\begin{analysisbox}{第 137 题}{答案：D}

\inlineanswer{D}

\analysiscontent{工作M有紧后O和P。O:最迟完成20天;P:最迟完成12天。M最迟完成=min(20-0,12-0)=12天。但需减去持续时间...M最迟完成=min(20,12)=12天,但O持续13天,最迟开始=20-13=7天;P持续8天,最迟开始=12-8=4天。M最迟完成=min(7,4)=4天(B)。}

\end{analysisbox}

\begin{analysisbox}{第 138 题}{答案：C}

\inlineanswer{C}

\analysiscontent{双代号时标网络计划:时间坐标方向垂直向下(A错);节点中心必须对准时标位置(B正确);不能用水平虚箭线表示虚工作(C错);时间坐标可以是日历坐标(D错)。}

\end{analysisbox}

\begin{analysisbox}{第 139 题}{答案：D}

\inlineanswer{D}

\analysiscontent{双代号早时标网络计划:A错:不能直接显示最迟时间;B错:逻辑关系表达不清;C错:有虚箭线;D正确:自由时差可通过波形线长度比较。}

\end{analysisbox}

\begin{analysisbox}{第 140 题}{答案：B}

\inlineanswer{B}

\analysiscontent{线路中各项工作持续时间之和为线路长度(C正确)。A错:长度最长的线路为关键线路;B正确:可能有一条或多条关键线路;D错:节点编号不必连续。}

\end{analysisbox}

\begin{analysisbox}{第 141 题}{答案：BD}

\inlineanswer{BD}

\analysiscontent{工艺关系是同一施工过程在不同施工段上的关系,还是不同施工过程在同一施工段上的关系?工艺关系是施工工艺决定的先后关系。A1B1属于工艺关系。答案B(A2B2)是同一施工段上不同过程的工艺关系。}

\end{analysisbox}

\begin{analysisbox}{第 142 题}{答案：AB}

\inlineanswer{AB}

\analysiscontent{根据双代号网络图绘图规则:节点编号有误(A)、存在循环回路(B)、多个起点节点(C)。D错:应为多个起点节点(非终点)。答案C:多个起点节点。}

\end{analysisbox}

\begin{analysisbox}{第 143 题}{答案：BCD}

\inlineanswer{BCD}

\analysiscontent{根据双代号网络图,工作I的最早开始时间需要从起点到I的各条线路中取最大值。答案D:第16天。}

\end{analysisbox}

\begin{analysisbox}{第 144 题}{答案：A}

\inlineanswer{A}

\analysiscontent{双代号时标网络计划中,工作F和H的最迟完成时间需要通过计算确定。答案C:第7天、第11天。}

\end{analysisbox}

\begin{analysisbox}{第 145 题}{答案：BD}

\inlineanswer{BD}

\analysiscontent{网络计划中工作D的总时差和自由时差:总时差=最迟完成-最早完成=2天;自由时差=紧后最早开始-最早完成=2天。答案D:2和2。}

\end{analysisbox}

\end{document}
